\documentclass[11pt]{article}
\usepackage[utf8]{inputenc}
\usepackage{amsmath,amssymb,amsfonts,amsthm}
\usepackage{geometry}
\usepackage{booktabs}
\usepackage{array}
\usepackage{graphicx}
\usepackage{verbatim}

\geometry{a4paper, margin=1in}

\title{CS282 Machine Learning: Homework 1 Solutions}
\author{Jianxuan Yao}
\date{2026-10-25}

\begin{document}

\maketitle

\section*{Problem 1}

\subsection*{(i)}

Re-arrange the $N$ data points in ascending order of $q(x_i)$. Then, they can be divided into 3 conditions:

\begin{enumerate}
    \item $-1, \dots, -1, +1, \dots, +1, -1, \dots, -1$
    \item $+1, \dots, +1$
    \item $-1, \dots, -1$
\end{enumerate}

It can be convert to a bars and stars problem.

The number of the first and the second condition is $\binom{N+1}{2}$, and the number of the third condition is $1$.

Therefore, the total number of different labelings is $m_\mathcal{H}(N) = \binom{N+1}{2} + 1$.

\subsection*{(ii)}

$$
d_{\text{vc}}(\mathcal{H}) = 2
$$

The smallest breakpoint is $3$.

Let $1 < q(x_1) < q(x_2)$.

We can construct $a$ and $b$ to shatter the two points:

\begin{enumerate}
    \item $a = \frac{q(x_1) + q(x_2)}{2}$, $b = \frac{q(x_1) + q(x_2)}{2}$: -1, -1
    \item $a = \frac{q(x_1) + q(x_2)}{2}$, $b = q(x_2) + 1$: -1, +1
    \item $a = q(x_1) - 1$, $b = \frac{q(x_1) + q(x_2)}{2}$: +1, -1
    \item $a = q(x_1) - 1$, $b = q(x_2) + 1$: +1, +1
\end{enumerate}

And for three points $q(x_1) < q(x_2) < q(x_3)$, we cannot construct parameters to let them be $+1, -1, +1$, because when $a < q(x_1) < b$ and $a < q(x_3) < b$, we must have $a < q(x_2) < b$, which means $h_{a,b}(x_2) = +1$.

\subsection*{(iii)}

\begin{enumerate}
    \item[1. ]
    
    $$
    4 m_\mathcal{H}(2N) \exp\left(-\frac{N\epsilon^2}{8}\right) = \delta
    $$

    $$
    \epsilon_\delta = \sqrt{\frac{8}{N} \ln\left(\frac{4 m_\mathcal{H}(2N)}{\delta}\right)}
    $$

    \item[2. ]
    
    $$
    m_\mathcal{H}(2N) = \binom{2N+1}{2} + 1 = 2N^2 + N + 1 = 2001001
    $$

    $$
    \frac{4 m_\mathcal{H}(2N)}{\delta} = \frac{4 \times 2,001,001}{0.05} = 160,080,080
    $$

    $$
    \ln(160,080,080) \approx 18.891185
    $$

    $$
    \epsilon_\delta \approx 0.389
    $$

    \item[3. ]
    
    A larger $\epsilon_\delta$ does not necessarily imply a poorer actual generalization performance.

    The VC generalization bound is a theoretical upper bound derived for the worst-case scenario. In practical applications, data often exhibits favorable geometric characteristics such as clustering, low-dimensional manifold structure, or high signal-to-noise ratio. Consequently, the actual generalization error $\vert E_{\text{out}}(g) - E_{\text{in}}(g) \vert$ is often much smaller than this theoretical upper bound $\epsilon_\delta$.
\end{enumerate}

\newpage

\section*{Problem 2}

\subsection*{(i)}

$$
\begin{aligned}
E_{D, \epsilon} [(g^{(D)}(x^\star) - y^\star)^2]
&= E_{D, \epsilon} [((g^{(D)}(x^\star) - \bar{g}(x^\star)) + (\bar{g}(x^\star) - f(x^\star)) - \epsilon)^2] \\
&=\quad  E_{D, \epsilon}((g^{(D)}(x^\star) - \bar{g}(x^\star))^2) \\
&\quad + E_{D, \epsilon}((\bar{g}(x^\star) - f(x^\star))^2) \\
&\quad + E_{D, \epsilon}(\epsilon^2) \\
&\quad + 2 E_{D, \epsilon}((g^{(D)}(x^\star) - \bar{g}(x^\star))(\bar{g}(x^\star) - f(x^\star))) \\
&\quad + 2 E_{D, \epsilon}((g^{(D)}(x^\star) - \bar{g}(x^\star))(-\epsilon)) \\
&\quad + 2 E_{D, \epsilon}((\bar{g}(x^\star) - f(x^\star))(-\epsilon)) \\
&= var(x^\star) + bias^2(x^\star) + \sigma^2
\end{aligned}
$$

In the first cross term, $\bar{g}(x^\star) - f(x^\star)$ is fixed and $E_D(g^{(D)}(x^\star) - \bar{g}(x^\star)) = E_{D}(g^{(D)}(x^\star)) - \bar{g}(x^\star) = 0$.

In the second and third cross term, $E_\epsilon(\epsilon) = 0$.

The bias measures the difference between the expected prediction of the model and the true function value. Always it is corresponding to underfitting.

The variance measures the variability of the model's predictions for different training datasets. Always it is corresponding to overfitting.

The noise term represents the irreducible error in the data, which cannot be eliminated by any model.

\subsection*{(ii)}

\begin{enumerate}

    \item[1. ]
    
    $$
    \Vert y - Xw \Vert _2^2 = (y - Xw)^\top (y - Xw) = y^\top y - 2w^\top X^\top y + w^\top X^\top X w
    $$

    $$
    w^\top X^\top X w = w^\top I_d w = \Vert w \Vert _2^2 = \sum_{j=1}^d w_j^2
    $$

    $X^\top y
    =
    \begin{bmatrix}
        X_{\cdot 1}^\top y \\
        \vdots \\
        X_{\cdot d}^\top y
    \end{bmatrix}$
    , so

    $$
    w^\top X^\top y = \sum_{j=1}^d w_j (X_{\cdot j}^\top y)
    $$

    $$
    \lambda \Vert w \Vert_1 = \lambda \sum_{j=1}^d \vert w_j \vert
    $$

    $$
    \begin{aligned}
    J_{\text{L}}(w)
    &= \frac{1}{2}\Vert y \Vert _2^2 - \sum_{j=1}^d w_j (X_{\cdot j}^\top y) + \frac{1}{2}\sum_{j=1}^d w_j^2 + \lambda \sum_{j=1}^d \vert w_j \vert \\
    &= \frac{1}{2}\Vert y \Vert _2^2 + \sum_{j=1}^d \left( \frac{1}{2} w_j^2 - w_j (X_{\cdot j}^\top y) + \lambda \vert w_j \vert \right)
    \end{aligned}
    $$

    $$
    C(y) = \frac{1}{2}\Vert y \Vert _2^2
    ,\quad
    \phi(X_{\cdot j}, y, w_j, \lambda) = \frac{1}{2} w_j^2 - (X_{\cdot j}^\top y) w_j + \lambda \vert w_j \vert
    $$

    $X^\top X = I_d$, so the expansion of $w^\top X^\top X w$ does not contain any cross terms $w_j w_k$ for $j \neq k$. Therefore, the objective function can be decomposed into a sum of functions, each depending only on a single weight $w_j$.

    \item[2. ]

    When $w_{\text{L}, j}^* > 0$:

    $$
    \phi(w_j) = \frac{1}{2}w_j^2 - (X_{\cdot j}^\top y) w_j + \lambda w_j
    $$

    $$
    \frac{\partial \phi}{\partial w_j} = w_j - X_{\cdot j}^\top y + \lambda = 0
    $$

    $$
    w_j = X_{\cdot j}^\top y - \lambda
    $$

    $$
    w_{\text{L}, j}^* = X_{\cdot j}^\top y - \lambda
    $$

    \item[3. ]

    When $w_{\text{L}, j}^* < 0$:

    $$
    \phi(w_j) = \frac{1}{2}w_j^2 - (X_{\cdot j}^\top y) w_j - \lambda w_j
    $$

    $$
    \frac{\partial \phi}{\partial w_j} = w_j - X_{\cdot j}^\top y - \lambda = 0
    $$

    $$
    w_{\text{L}, j}^* = X_{\cdot j}^\top y + \lambda
    $$

    \item[4. ]

    When $w_{\text{L}, j}^* = 0$:

    $$
    \partial \phi(0) = -X_{\cdot j}^\top y + \lambda [-1, 1] = \left[-X_{\cdot j}^\top y - \lambda, -X_{\cdot j}^\top y + \lambda\right]
    $$

    $0 \in \left[-X_{\cdot j}^\top y - \lambda, -X_{\cdot j}^\top y + \lambda\right]$, so $\vert X_{\cdot j}^\top y \vert \le \lambda$

    $$
    w_{\text{L}, j}^* = 0 \Leftrightarrow \vert X_{\cdot j}^\top y \vert \le \lambda
    $$

    Based on the above three cases,
    
    $$
    w_{\text{L}, j}^*
    = \mathcal{S}_\lambda(X_{\cdot j}^\top y)
    = \text{sign}(X_{\cdot j}^\top y) \cdot \max\left(\vert X_{\cdot j}^\top y \vert - \lambda, 0\right)
    $$

    If $\vert X_{\cdot j}^\top y \vert \leq \lambda$, then the corresponding weight of the feature will be $w_{\text{L}, j}^* = 0$.

    What's more, this phenomenon can be interpreted as a form of L1 regularization, which owns a geometric angle.
\end{enumerate}

\subsection*{(iii)}

$$
J_{\text{R}}(w) = \frac{1}{2}\Vert y - Xw \Vert _2^2 + \frac{\lambda}{2}\Vert w \Vert _2^2
$$

$$
\nabla_w J_{\text{R}}(w) = -X^\top (y - Xw) + \lambda w = 0
$$

$$
w_{\text{R}}^* = \frac{1}{1 + \lambda} X^\top y
$$

$$
w_{\text{R}, j}^* = 0 \Leftrightarrow X_{\cdot j}^\top y = 0
$$

Ridge regression does not have the same sparsity-inducing property as Lasso regression.

In Ridge regression, if and only if the feature $X_{\cdot j}$ is orthogonal to the target vector $y$, the corresponding weight $w_{\text{R}, j}^*$ will be zero. Otherwise, Ridge regression will always assign a non-zero weight to each feature.

In LASSO regression, $\vert X_{\cdot j}^\top y \vert \le \lambda$ is an interval with positive measure, so it has a innate sparsity generating property.

\subsection*{(iv)}

$$
w_{\text{L}}^* = (1.3, 0, 0, 2.1)^\top
$$

$$
w_{\text{R}}^*
= \left(1.2, -\frac{4}{15}, \frac{1}{15}, -\frac{26}{15}\right)^\top
$$

LASSO regression has 2 zero weights, while Ridge regression has no zero weights.

Since the original correlation strengths of the second and third terms, $\vert z_2\vert  = 0.4$ and $\vert z_3\vert  = 0.1$, do not exceed the penalty threshold $\lambda = 0.5$, LASSO precisely sets them to 0, resulting in sparsity; whereas Ridge regression outputs non-zero solutions for all non-zero components and cannot eliminate unimportant features.

Ridge regression shrinks every weight by a factor of $\frac{1}{1 + \lambda}$, while LASSO regression shrinks weights larger than $\lambda$ by a constant $\lambda$ and set weights smaller than $\lambda$ to zero.

\newpage

\section*{Problem 3}

\subsection*{(i)}

$$
\begin{aligned}
J_\lambda(\theta_1, \theta_0)
&= \frac{1}{6} \sum_{n=1}^3 (\theta_1 x_n + \theta_0 - y_n)^2 + \frac{\lambda}{2}\theta_1^2 \\
&= \frac{1}{6} \left[ (-\theta_1 + \theta_0 + 1)^2 + (\theta_1 + \theta_0 - 3)^2 + (2\theta_1 + \theta_0 - 5)^2 \right] + \frac{\lambda}{2}\theta_1^2 \\
&= \frac{1}{6}\left[ 3\theta_0^2 + (6 + 3\lambda)\theta_1^2 + 4\theta_0\theta_1 - 14\theta_0 - 28\theta_1 + 35 \right]
\end{aligned}
$$

$$
\frac{\partial J_\lambda}{\partial \theta_1} = (2 + \lambda)\theta_1 + \frac{2}{3}\theta_0 - \frac{14}{3}
$$

$$
\frac{\partial J_\lambda}{\partial \theta_0} = \frac{2}{3}\theta_1 + \theta_0 - \frac{7}{3}
$$

$$
\nabla J_\lambda(\mathbf{u})
=
\begin{pmatrix}
    \frac{\partial J_\lambda}{\partial \theta_1} \\
    \frac{\partial J_\lambda}{\partial \theta_0}
\end{pmatrix}
=
\begin{pmatrix}
    (2 + \lambda)\theta_1 + \frac{2}{3}\theta_0 - \frac{14}{3} \\
    \frac{2}{3}\theta_1 + \theta_0 - \frac{7}{3}
\end{pmatrix}
$$

$$
\frac{\partial^2 J_\lambda}{\partial \theta_1^2} = 2 + \lambda
$$

$$
\frac{\partial^2 J_\lambda}{\partial \theta_1 \partial \theta_0}
= \frac{\partial^2 J_\lambda}{\partial \theta_0 \partial \theta_1}
= \frac{2}{3}
$$

$$
\frac{\partial^2 J_\lambda}{\partial \theta_0^2} = 1
$$

$$
H_\lambda
=
\begin{pmatrix} 
    2 + \lambda & \frac{2}{3} \\
    \frac{2}{3} & 1
\end{pmatrix}
$$

$$
\nabla J_\lambda(\mathbf{u}^\star) = \mathbf{0}
$$

$$
\theta_1 = \frac{28}{14 + 9\lambda}
,\quad
\theta_0 = \frac{7}{3} - \frac{2}{3}\left(\frac{28}{14 + 9\lambda}\right)
$$

When $\lambda = 0$, $
u^\star
=
\begin{pmatrix}
    2 \\
    1
\end{pmatrix}
$;
When $\lambda = 2$, $u^\star
=
\begin{pmatrix} 
    \frac{7}{8} \\
    \frac{7}{4}
\end{pmatrix}$.

\subsection*{(ii)}

According to the centre computed in (i), $\lambda = 0$ corresponds to B, D, F, and $\lambda = 2$ corresponds to A, C, E.

$$
\nabla J_\lambda(0, 0)
=
\begin{pmatrix}
    -\frac{14}{3} \\
    -\frac{7}{3}
\end{pmatrix}
$$

$$
\mathbf{u}^{(1)}
=
\alpha
\begin{pmatrix}
    \frac{14}{3} \\
    \frac{7}{3}
\end{pmatrix}
$$

$\alpha = 0.1$, $u_1
\approx
\begin{pmatrix}
    0.4667 \\
    0.2333
\end{pmatrix}$;

$\alpha = 0.4$, $u_1
\approx
\begin{pmatrix}
    1.8667 \\
    0.9333
\end{pmatrix}$;

$\alpha = 0.6$, $u_1
=
\begin{pmatrix}
    2.8 \\
    1.4
\end{pmatrix}$.

As a result:

A: $\alpha = 0.4, \ \lambda = 2$

B: $\alpha = 0.1, \ \lambda = 0$

C: $\alpha = 0.6, \ \lambda = 2$

D: $\alpha = 0.6, \ \lambda = 0$

E: $\alpha = 0.1, \ \lambda = 2$

F: $\alpha = 0.4, \ \lambda = 0$

\subsection*{(iii)}

\begin{enumerate}
    \item[1. ]
    
    $$
    f(\mathbf{u}) \le f(\mathbf{v}) + \nabla f(\mathbf{v})^\top (\mathbf{u} - \mathbf{v}) + \frac{L}{2}\Vert  \mathbf{u} - \mathbf{v}\Vert  _2^2
    $$

    $$
    f(\mathbf{x}^+) \le f(\mathbf{x}) + \nabla f(\mathbf{x})^\top (\mathbf{x}^+ - \mathbf{x}) + \frac{L}{2}\Vert  \mathbf{x}^+ - \mathbf{x}\Vert _2^2
    $$

    $\mathbf{x}^+ - \mathbf{x} = -\frac{1}{L}\nabla f(\mathbf{x})$, so

    $$
    \nabla f(\mathbf{x})^\top (\mathbf{x}^+ - \mathbf{x}) = \nabla f(\mathbf{x})^\top \left(-\frac{1}{L}\nabla f(\mathbf{x})\right) = -\frac{1}{L}\Vert \nabla f(\mathbf{x})\Vert _2^2
    $$

    $$
    \frac{L}{2}\Vert \mathbf{x}^+ - \mathbf{x}\Vert _2^2 = \frac{L}{2}\left\Vert -\frac{1}{L}\nabla f(\mathbf{x})\right\Vert _2^2 = \frac{L}{2} \cdot \frac{1}{L^2}\Vert \nabla f(\mathbf{x})\Vert _2^2 = \frac{1}{2L}\Vert \nabla f(\mathbf{x})\Vert _2^2
    $$

    $$
    f(\mathbf{x}^+) \le f(\mathbf{x}) - \frac{1}{L}\Vert \nabla f(\mathbf{x})\Vert _2^2 + \frac{1}{2L}\Vert \nabla f(\mathbf{x})\Vert _2^2 = f(\mathbf{x}) - \frac{1}{2L}\Vert \nabla f(\mathbf{x})\Vert _2^2
    $$

    \item[2. ]
    
    $$
    f(\mathbf{x})
    \le
    f(\mathbf{x}^\star) + \nabla f(\mathbf{x})^\top (\mathbf{x} - \mathbf{x}^\star)
    $$

    $$
    f(\mathbf{x}^+)
    \le
    f(\mathbf{x}^\star) + \nabla f(\mathbf{x})^\top (\mathbf{x} - \mathbf{x}^\star) - \frac{1}{2L}\Vert \nabla f(\mathbf{x})\Vert _2^2
    $$

    $$
    f(\mathbf{x}^+) - f(\mathbf{x}^\star)
    \le \nabla f(\mathbf{x})^\top (\mathbf{x} - \mathbf{x}^\star) - \frac{1}{2L}\Vert \nabla f(\mathbf{x})\Vert _2^2
    = \frac{L}{2} \left[ \frac{2}{L}\nabla f(\mathbf{x})^\top (\mathbf{x} - \mathbf{x}^\star) - \frac{1}{L^2}\Vert \nabla f(\mathbf{x})\Vert _2^2 \right]
    $$

    $$
    \frac{2}{L}\nabla f(\mathbf{x})^\top (\mathbf{x} - \mathbf{x}^\star) - \frac{1}{L^2}\Vert \nabla f(\mathbf{x})\Vert _2^2
    = 2(\mathbf{x} - \mathbf{x}^+)^\top (\mathbf{x} - \mathbf{x}^\star) - \Vert \mathbf{x} - \mathbf{x}^+\Vert _2^2
    $$

    Let $\mathbf{a} = \mathbf{x}^+ - \mathbf{x}^\star$ and $\mathbf{b} = \mathbf{x} - \mathbf{x}^+$, we have

    $$
    \mathbf{a} + \mathbf{b} = (\mathbf{x}^+ - \mathbf{x}^\star) + (\mathbf{x} - \mathbf{x}^+) = \mathbf{x} - \mathbf{x}^\star
    $$

    According to $2\mathbf{a}^\top \mathbf{b} + \Vert \mathbf{b}\Vert _2^2 = \Vert \mathbf{a} + \mathbf{b}\Vert _2^2 - \Vert \mathbf{a}\Vert _2^2$, 

    $$
    \begin{aligned}
        \Vert \mathbf{x} - \mathbf{x}^\star\Vert _2^2 - \Vert \mathbf{x}^+ - \mathbf{x}^\star\Vert _2^2
        &= \Vert \mathbf{a} + \mathbf{b}\Vert _2^2 - \Vert \mathbf{a}\Vert _2^2 \\
        &= 2\mathbf{a}^\top \mathbf{b} + \Vert \mathbf{b}\Vert _2^2 \\
        &= 2(\mathbf{x}^+ - \mathbf{x}^\star)^\top (\mathbf{x} - \mathbf{x}^+) + \Vert \mathbf{x} - \mathbf{x}^+\Vert _2^2
    \end{aligned}
    $$

    $$
    \begin{aligned}
        2(\mathbf{x}^+ - \mathbf{x}^\star)^\top (\mathbf{x} - \mathbf{x}^+) + \Vert \mathbf{x} - \mathbf{x}^+\Vert _2^2
        &= 2\left[(\mathbf{x} - \mathbf{x}^\star) - (\mathbf{x} - \mathbf{x}^+)\right]^\top (\mathbf{x} - \mathbf{x}^+) + \Vert \mathbf{x} - \mathbf{x}^+\Vert _2^2 \\
        &= 2(\mathbf{x} - \mathbf{x}^\star)^\top (\mathbf{x} - \mathbf{x}^+) - 2\Vert \mathbf{x} - \mathbf{x}^+\Vert _2^2 + \Vert \mathbf{x} - \mathbf{x}^+\Vert _2^2 \\
        &= 2(\mathbf{x} - \mathbf{x}^\star)^\top (\mathbf{x} - \mathbf{x}^+) - \Vert \mathbf{x} - \mathbf{x}^+\Vert _2^2
    \end{aligned}
    $$

    $$
    \frac{2}{L}\nabla f(\mathbf{x})^\top (\mathbf{x} - \mathbf{x}^\star) - \frac{1}{L^2}\Vert \nabla f(\mathbf{x})\Vert _2^2
    =
    \Vert \mathbf{x} - \mathbf{x}^\star\Vert _2^2 - \Vert \mathbf{x}^+ - \mathbf{x}^\star\Vert _2^2
    $$

    $$
    f(\mathbf{x}^+) - f(\mathbf{x}^\star)
    \le
    \frac{L}{2} \left( \Vert \mathbf{x} - \mathbf{x}^\star\Vert _2^2 - \Vert \mathbf{x}^+ - \mathbf{x}^\star\Vert _2^2 \right)
    $$
\end{enumerate}

\newpage

\section*{Problem 4}

\subsection*{(i)}

\begin{enumerate}
    \item[1. ]
    
    No. The decision boundary of an affine linear classifier is a point. However, the data needs at least two points to be separated.

    \item[2. ]
    
    $z_1 = (-2, 4)^\top$, $y_1 = +1$;

    $z_2 = (-1, 1)^\top$, $y_2 = -1$;

    $z_3 = (0, 0)^\top$, $y_3 = -1$;

    $z_4 = (1, 1)^\top$, $y_4 = -1$;

    $z_5 = (2, 4)^\top$, $y_5 = +1$.

    Every point is on a parabola in a two-dimensional plane.

    They are dividable. We can construct $z_2 = 2$ to separate them.

    \item[3. ]
    
    $\phi(x) = \begin{pmatrix} x \\ x^2 \end{pmatrix}$ and $\phi(x') = \begin{pmatrix} x' \\ (x')^2 \end{pmatrix}$, so

    $$
    K(x, x') = x x' + x^2 (x')^2 = x x' + (x x')^2
    $$
\end{enumerate}

\subsection*{(ii)}

\begin{enumerate}
    \item[1. ]
    
    The samples are symmetric about $z_1$ axis, so the decision hyperplane should be parallel to $z_1$ axis.

    $$
    \begin{cases}
        w_2 \cdot 4 + b = 1 \\
        w_2 \cdot 1 + b = -1
    \end{cases}
    $$

    $$
    w^\star
    =
    \begin{pmatrix}
        0 \\
        \frac{2}{3}
    \end{pmatrix}
    , \quad
    b^\star = -\frac{5}{3}
    $$

    We need to find a group of $\alpha$ to satisfy the KKT condition:

    $$
    \begin{cases}
        \mathbf{w} &= \sum_{i=1}^5 \alpha_i y_i \mathbf{z}_i \\
        \sum_{i=1}^5 \alpha_i y_i &= 0 \\
        y_i(\mathbf{w}^\top \mathbf{z}_i + b) &\ge 1, \quad \forall i=1,\dots,5 \\
        \alpha_i &\ge 0, \quad \forall i=1,\dots,5 \\
        \alpha_i \left[ y_i(\mathbf{w}^\top \mathbf{z}_i + b) - 1 \right] &= 0, \quad \forall i=1,\dots,5
    \end{cases}
    $$

    Let $\alpha_1 = \alpha_5 = \frac{1}{9}$, $\alpha_2 = \alpha_4 = \frac{1}{9}$, $\alpha_3 = 0$. Eazy to find they satisfy the KKT condition.

    \item[2. ]
    
    $$
    \gamma = \frac{1}{\Vert \mathbf{w}^\star\Vert _2} = \frac{3}{2}
    $$

    $$
    M = 3
    $$

    $$
    \mathbf{w}^{\star\top} \mathbf{z} + b^\star = +1 \iff z_2 = 4
    $$

    $$
    \mathbf{w}^{\star\top} \mathbf{z} + b^\star = 0 \iff z_2 = \frac{5}{2}
    $$

    $$
    \mathbf{w}^{\star\top} \mathbf{z} + b^\star = -1 \iff z_2 = 1
    $$

    Sample $1, 2, 4, 5$ satisfy $y_i((\mathbf{w}^\star)^\top \phi(x_i) + b^\star) = 1$.

    \item[3. ]

    The original decision boundary is:
    
    $$
    x = \pm \frac{\sqrt{10}}{2}
    $$

    $\left(-\infty, -\sqrt{\frac{5}{2}}\,\right] \cup \left[\sqrt{\frac{5}{2}}, +\infty\right)$ for $+1$;

    $\left(-\sqrt{\frac{5}{2}}, \sqrt{\frac{5}{2}}\right)$ for $-1$.
\end{enumerate}

\subsection*{(iii)}

\begin{enumerate}
    \item[1. ] $\xi_i = 0$
    
    $$
    \xi_i = 0 \iff 1 - y_i f(\mathbf{z}_i) \le 0 \iff y_i f(\mathbf{z}_i) \ge 1
    $$

    The sample is one the margin boundary of its class or in the correct side.

    Correctly classified.

    \item[2. ] $0 < \xi_i \leq 1$
    
    Condition $0 < \xi_i < 1$:

    $$
    0 < 1 - \xi_i < 1 \implies 0 < y_i f(\mathbf{z}_i) < 1
    $$

    Bewteen the margin boundary and the decision boundary of its class.

    Correctly classified.

    Condition $\xi_i = 1$:

    $$
    y_i f(\mathbf{z}_i) = 1 - 1 = 0 \implies f(\mathbf{z}_i) = 0
    $$

    On the decision boundary.

    If $y_i = +1$, it classified correctly; if $y_i = -1$, it classified incorrectly.

    \item[3. ] $\xi_i > 1$
    
    $$
    y_i f(\mathbf{z}_i) = 1 - \xi_i < 0
    $$

    On the wrong side of the decision boundary of its class.

    Misclassified.
\end{enumerate}

\subsection*{(iv)}

\begin{enumerate}
    \item[1. ]
    $$
    \begin{aligned}
        \Vert \psi(\mathbf{u}) - \psi(\mathbf{v})\Vert ^2
        &= \langle \psi(\mathbf{u}) - \psi(\mathbf{v}), \, \psi(\mathbf{u}) - \psi(\mathbf{v}) \rangle \\
        &= \langle \psi(\mathbf{u}), \psi(\mathbf{u}) \rangle - 2 \langle \psi(\mathbf{u}), \psi(\mathbf{v}) \rangle + \langle \psi(\mathbf{v}), \psi(\mathbf{v}) \rangle
    \end{aligned}
    $$

    $$
    \Vert \psi(\mathbf{u}) - \psi(\mathbf{v})\Vert ^2 = K(\mathbf{u}, \mathbf{u}) - 2K(\mathbf{u}, \mathbf{v}) + K(\mathbf{v}, \mathbf{v})
    $$

    \item[2. ]

    $$
    K(\mathbf{u}, \mathbf{u})
    = \exp\left(-\frac{\Vert \mathbf{u} - \mathbf{u}\Vert _2^2}{2\sigma^2}\right)
    = \exp(0)
    = 1
    $$

    $$
    \Vert \psi(\mathbf{u})\Vert 
    = \sqrt{\langle \psi(\mathbf{u}), \psi(\mathbf{u}) \rangle}
    = \sqrt{K(\mathbf{u}, \mathbf{u})}
    = 1
    $$

    $$
    \begin{aligned}
        \Vert \psi(\mathbf{u}) - \psi(\mathbf{v})\Vert ^2
        &= K(\mathbf{u}, \mathbf{u}) - 2K(\mathbf{u}, \mathbf{v}) + K(\mathbf{v}, \mathbf{v}) \\
        &= 1 - 2\exp\left(-\frac{\Vert \mathbf{u} - \mathbf{v}\Vert _2^2}{2\sigma^2}\right) + 1 \\
        &= 2 - 2\exp\left(-\frac{\Vert \mathbf{u} - \mathbf{v}\Vert _2^2}{2\sigma^2}\right)
    \end{aligned}
    $$

    When $\Vert \mathbf{u} - \mathbf{v}\Vert _2 \to \infty$, 

    $$
    \lim_{\Vert \mathbf{u} - \mathbf{v}\Vert _2 \to \infty} \exp\left(-\frac{\Vert \mathbf{u} - \mathbf{v}\Vert _2^2}{2\sigma^2}\right) = 0
    $$

    $$
    \lim_{\Vert \mathbf{u} - \mathbf{v}\Vert _2 \to \infty} \Vert \psi(\mathbf{u}) - \psi(\mathbf{v})\Vert  = \sqrt{2}
    $$
\end{enumerate}

\subsection*{(v)}

\begin{enumerate}
    \item[1. ] Must be a valid kernel function.
    
    Define $\phi(\mathbf{x}) = \begin{pmatrix} \phi_1(\mathbf{x}) \\ \phi_2(\mathbf{x}) \end{pmatrix}$.

    $$
    \begin{aligned}
        \langle \phi(\mathbf{u}), \phi(\mathbf{v}) \rangle 
        &= \left\langle \begin{pmatrix} \phi_1(\mathbf{u}) \\ 
        \phi_2(\mathbf{u}) \end{pmatrix}, \, \begin{pmatrix} \phi_1(\mathbf{v}) \\ \phi_2(\mathbf{v}) \end{pmatrix} \right\rangle \\
        &= \langle \phi_1(\mathbf{u}), \phi_1(\mathbf{v}) \rangle + \langle \phi_2(\mathbf{u}), \phi_2(\mathbf{v}) \rangle \\
        &= K_1(\mathbf{u}, \mathbf{v}) + K_2(\mathbf{u}, \mathbf{v}) \\
        &= K(\mathbf{u}, \mathbf{v})
    \end{aligned}
    $$

    \item[2. ] May not be a valid kernel function.
    
    Let $K_1(u, v) = u v$ and $K_2(u, v) = 2 u v$.

    $$
    K(u, v) = K_1(u, v) - K_2(u, v) = uv - 2uv = -uv
    $$

    $$
    K(1, 1) = -1 \cdot 1 = -1 < 0
    $$

    It contradicts the property of a valid kernel function, which requires $K(x, x) \ge 0$ for all $x$.
\end{enumerate}

\newpage

\section*{Problem 5}

\subsection*{(i)}

\begin{enumerate}
    \item[1. ]

    $$
    \phi_p(x) = (1, x, \dots, x^p)^\top
    ,\quad
    \Phi = \begin{pmatrix}
        1 & x_1 & \cdots & x_1^p \\
        1 & x_2 & \cdots & x_2^p \\
        \vdots & \vdots & \ddots & \vdots \\
        1 & x_n & \cdots & x_n^p
    \end{pmatrix}
    $$

    $$
    (g_{p,w}(x_1), \dots, g_{p,w}(x_n))^\top = \Phi w
    $$

    \item[2. ]

    Let $R = \operatorname{diag}(0, 1, \dots, 1) \in \mathbb{R}^{(p+1)\times(p+1)}$.

    $$
    J_{S,p,\lambda}(w) = \frac{1}{2n}\Vert \Phi w - y\Vert_2^2 + \frac{\lambda}{2}w^\top Rw
    $$

    $$
    \nabla_w J_{S,p,\lambda}(w)
    = \frac{1}{n}\Phi^\top(\Phi w-y) + \lambda Rw = 0
    $$

    $$
    (\Phi^\top\Phi + n\lambda R)w = \Phi^\top y
    $$

    $$
    w^\star = (\Phi^\top\Phi + n\lambda R)^{-1}\Phi^\top y
    $$

    Since $\Phi$ has full column rank, for any $v \neq 0$,

    $$
    v^\top(\Phi^\top\Phi + n\lambda R)v
    = \Vert\Phi v\Vert_2^2 + n\lambda\sum_{j=1}^p v_j^2 > 0.
    $$

    Therefore, $w^\star$ is the unique minimizer.

    \item[3. ]

    $R^\top R = R$, so

    $$
    w^\star = \arg\min_w
    \left\Vert
    \begin{pmatrix} \Phi \\ \sqrt{n\lambda}R \end{pmatrix}w
    - \begin{pmatrix} y \\ 0 \end{pmatrix}
    \right\Vert_2^2.
    $$

    Use \texttt{np.linalg.lstsq} on this augmented system to avoid explicitly forming $\Phi^\top\Phi$. The implementation is included below.
\end{enumerate}

\subsection*{(ii)}

\begin{enumerate}
    \item[1. ]

    \begin{center}
    \begin{tabular}[t]{ccc}
        \toprule
        $p$ & Training MSE & Validation MSE \\
        \midrule
        1 & 0.23540310 & 0.23144095 \\
        2 & 0.20670953 & 0.26599808 \\
        3 & 0.03322276 & 0.04401319 \\
        4 & 0.03291933 & 0.04680137 \\
        5 & 0.03285123 & 0.04307267 \\
        6 & 0.03263838 & 0.04063874 \\
        7 & 0.03067714 & 0.05759251 \\
        8 & 0.02894150 & 0.06978324 \\
        9 & 0.02785983 & 0.08036738 \\
        10 & 0.02753541 & 0.06886355 \\
        \bottomrule
    \end{tabular}
    \end{center}

    \begin{center}
        \includegraphics[width=0.9\linewidth]{problem5_figures/errors.pdf}
    \end{center}

    \item[2. ]

    $$
    p_{\text{train}}^\star = 10
    ,\quad
    p_{\text{val}}^\star = 6
    $$

    $$
    \mathcal{H}_p \subseteq \mathcal{H}_{p+1}
    \implies
    \min_{g\in\mathcal{H}_{p+1}} E_{\text{train}}(g)
    \le \min_{g\in\mathcal{H}_p} E_{\text{train}}(g).
    $$

    Higher degrees can fit more training noise. Training MSE decreases, while validation MSE need not decrease because of overfitting.

    \item[3. ]

    \begin{center}
        \raisebox{-\height}{\includegraphics[width=\linewidth]{problem5_figures/fits.pdf}}
    \end{center}
\end{enumerate}

\subsection*{(iii)}

\begin{enumerate}
    \item[1. ]

    $$
    S_{-j} = \bigcup_{k\neq j}D_k
    ,\quad |S_{-j}| = 32
    ,\quad |D_j| = 8
    $$

    $$
    w^{(-j)}_{p,\lambda}
    = \arg\min_w J_{S_{-j},p,\lambda}(w)
    $$

    $$
    E_{\text{CV}}(p,\lambda)
    = \frac{1}{5}\sum_{j=1}^5\frac{1}{8}
    \sum_{i\in D_j}\left(\phi_p(x_i)^\top w^{(-j)}_{p,\lambda}-y_i\right)^2
    $$

    \begin{center}
    \small
    \setlength{\tabcolsep}{5pt}
    \begin{tabular}{c r r r r r r}
        \toprule
        $p\backslash\lambda$ & $0$ & $10^{-4}$ & $10^{-3}$ & $10^{-2}$ & $10^{-1}$ & $1$ \\
        \midrule
        1 & 0.25817241 & 0.25816325 & 0.25808375 & 0.25755813 & 0.26894279 & 0.43655696 \\
        2 & 0.26527077 & 0.26520879 & 0.26466613 & 0.26049927 & 0.26438550 & 0.43502955 \\
        3 & 0.04155653 & 0.04158875 & 0.04327030 & 0.09618864 & 0.25590384 & 0.42935756 \\
        4 & 0.04332505 & 0.04302685 & 0.04390291 & 0.10274463 & 0.25918380 & 0.42837961 \\
        5 & 0.04321255 & 0.04442243 & 0.04962433 & 0.07502662 & 0.23558947 & 0.42914219 \\
        6 & 0.09179643 & 0.05288220 & 0.05122792 & 0.07655224 & 0.24035168 & 0.42882265 \\
        7 & 0.07899768 & 0.05151095 & 0.05483393 & 0.07317050 & 0.22180713 & 0.42911122 \\
        8 & 0.05719230 & 0.06815458 & 0.05962593 & 0.07327274 & 0.22596636 & 0.42910469 \\
        9 & 0.23272273 & 0.06535286 & 0.06093946 & 0.07515806 & 0.21337213 & 0.42890939 \\
        10 & 1.19941009 & 0.08338072 & 0.06789959 & 0.07513874 & 0.21671489 & 0.42902223 \\
        \bottomrule
    \end{tabular}
    \end{center}

    \item[2. ]

    $$
    (p^\star,\lambda^\star) = (3,0)
    ,\quad E_{\text{CV}}(3,0) = 0.04155653
    $$

    Compare $(E_{\text{CV}},p,\lambda)$ lexicographically using unrounded errors.

    \item[3. ]

    $$
    E_{\text{CV}}(10,0) = 1.19941009
    ,\quad
    E_{\text{CV}}(10,10^{-4}) = 0.08338072
    $$

    The unregularized degree-10 model is sensitive to the training fold. A small positive $\lambda$ penalizes large coefficients, reduces variance, and substantially improves validation performance.
\end{enumerate}

\subsection*{(iv)}

\begin{enumerate}
    \item[1. ]

    Refit $(p^\star,\lambda^\star) = (3,0)$ using all 40 examples:

    $$
    w^\star = \begin{pmatrix}
        0.0048801042 \\
        2.7215440602 \\
        -0.1052813046 \\
        -3.1209193461
    \end{pmatrix}
    $$

    $$
    g(x) = 0.0048801042 + 2.7215440602x - 0.1052813046x^2 - 3.1209193461x^3
    $$

    \item[2. ]

    Generate the 1000 test examples with seed $2026282$ after refitting.

    $$
    E_{\text{test}}(g) = \frac{1}{1000}\sum_{i=1}^{1000}(g(x_i)-y_i)^2 = 0.05507274
    $$

    \item[3. ]

    Each cross-validation model uses only 32 examples. Refitting uses all 40 examples to estimate the final parameters.

    If test error is used to tune $p$ or $\lambda$, the test set becomes part of model selection and its error can be optimistically biased. It should be used only for the final independent evaluation.
\end{enumerate}

\newpage

\subsection*{Code}

\texttt{problem5.py} (Python, NumPy, Matplotlib).

\begingroup
\footnotesize
\verbatiminput{problem5.py}
\endgroup

\end{document}